% from ega2-1-fr.tex
\oldpage[II]{86}

  et une section $s'$ de $\mathcal{L}'^{\otimes n}$ au-dessus de $X$ tels
  que $s'(x)\neq0$. Alors, si $\mathcal{L}$ est ample, il en est de même de
  $\mathcal{L}\otimes\mathcal{L}'$.
\end{enumerate}
\end{proposition}

\begin{proof}
La propriété (i) est conséquence évidente du critère a) de
(\hyperref[II.4.5.2-fr]{4.5.2}) puisque $X_{f^{\otimes n}}=X_f$. D'autre
part, si $\mathcal{L}$ est ample, pour tout $x\in X$ et tout voisinage $U$ de
$x$, il existe $m>0$ et $f\in\Gamma(X,\mathcal{L}^{\otimes m})$ tel que
$x\in X_f\subset U$ (\hyperref[II.4.5.2-fr]{4.5.2, a)})~; si
$f'\in\Gamma(X,\mathcal{L}'^{\otimes n})$ est tel que $f'(x)\neq0$, on aura
$s(x)\neq0$ pour
$s=f^{\otimes n}\otimes f'^{\otimes m}\in
\Gamma(X,(\mathcal{L}\otimes\mathcal{L}')^{\otimes mn})$, donc
$x\in X_s\subset X_f\subset U$, ce qui prouve que
$\mathcal{L}\otimes\mathcal{L}'$ est ample
(\hyperref[II.4.5.2-fr]{4.5.2, a)}).
\end{proof}

\begin{corollary}[4.5.7]
\label{II.4.5.7-fr}
Le produit tensoriel de deux $\mathcal{O}_X$-Modules amples est ample.
\end{corollary}

\begin{corollary}[4.5.8]
\label{II.4.5.8-fr}
Soient $\mathcal{L}$ un $\mathcal{O}_X$-Module ample, $\mathcal{L}'$ un
$\mathcal{O}_X$-Module inversible~; il existe alors un entier $n_0>0$ tel
que $\mathcal{L}^{\otimes n}\otimes\mathcal{L}'$ soit ample et engendré par
ses sections au-dessus de $X$ pour $n\geq n_0$.
\end{corollary}

\begin{proof}
En effet, il résulte de (\hyperref[II.4.5.5-fr]{4.5.5}) qu'il existe un entier
$m_0$ tel que $\mathcal{L}^{\otimes m}\otimes\mathcal{L}'$ soit engendré par
ses sections au-dessus de $X$ pour $m\geq m_0$~; d'après
(\hyperref[II.4.5.6-fr]{4.5.6}) on peut donc prendre $n_0=m_0+1$.
\end{proof}

\begin{remark}[4.5.9]
\label{II.4.5.9-fr}
Soit $P=\mathrm{H}^1(X,\mathcal{O}_X^*)$ le groupe des classes de
$\mathcal{O}_X$-Modules inversibles (\hyperref[0.5.4.7-fr]{0, 5.4.7}), et
soit $P^+$ la partie de $P$ formée des classes de faisceaux amples. Supposons
que $P^+$ soit non vide. Alors il résulte de
(\hyperref[II.4.5.7-fr]{4.5.7}) et
(\hyperref[II.4.5.8-fr]{4.5.8}) que l'on a
\[
  P^++P^+\subset P^+
  \qquad\text{et}\qquad
  P^+-P^+=P
\]
autrement dit $P^+\cup\{0\}$ est l'ensemble des éléments positifs dans $P$
pour une structure de préordre sur $P$ compatible avec sa structure de
groupe, et qui est même archimédienne en vertu de
(\hyperref[II.4.5.8-fr]{4.5.8}). C'est pourquoi on dit parfois «~faisceau
positif~» au lieu de faisceau ample et «~faisceau négatif~» pour l'inverse
d'un tel faisceau (terminologie que nous ne suivrons pas).
\end{remark}

\begin{proposition}[4.5.10]
\label{II.4.5.10-fr}
Soient $Y$ un schéma affine, $q:X\to Y$ un morphisme séparé quasi-compact,
$\mathcal{L}$ un $\mathcal{O}_X$-Module inversible.
\begin{enumerate}
  \item[\rm{(i)}] Si $\mathcal{L}$ est très ample relativement à $q$, alors
  $\mathcal{L}$ est ample.
  \item[\rm{(ii)}] Supposons en outre que le morphisme $q$ soit de type fini.
  Alors, pour que $\mathcal{L}$ soit ample, il faut et il suffit qu'il possède
  l'une des propriétés équivalentes suivantes~:
  \begin{enumerate}
    \item[\rm{e)}] Il existe $n_0>0$ tel que pour tout entier $n\geq n_0$,
    $\mathcal{L}^{\otimes n}$ soit très ample relativement à $q$.
    \item[\rm{e')}] Il existe $n>0$ tel que $\mathcal{L}^{\otimes n}$ soit
    très ample relativement à $q$.
  \end{enumerate}
\end{enumerate}
\end{proposition}

\begin{proof}
La première assertion résulte de la définition
(\hyperref[II.4.4.2-fr]{4.4.2}) d'un $\mathcal{O}_X$-Module très ample~: si
$A$ est l'anneau de $Y$, il y a un $A$-module $E$ et un homomorphisme
surjectif
\[
  \psi:q^*((\mathbf{S}(E))^{\sim})\longrightarrow
  \bigoplus_{n\geq0}\mathcal{L}^{\otimes n}
\]
tel que $i=r_{\mathcal{L},\psi}$ soit une immersion partout définie
$X\to P=\mathbf{P}(\widetilde{E})$ et que l'on ait
$\mathcal{L}=i^*(\mathcal{O}_P(1))$~; comme les $D_+(f)$ pour $f$ homogène
dans $(\mathbf{S}(E))_+$ forment une base de la topologie de $P$, et que
$i^{-1}(D_+(f))=X_{\psi^\flat(f)}$ par
(\hyperref[II.3.7.3.1-fr]{3.7.3.1}), on voit que la condition a) de
(\hyperref[II.4.5.2-fr]{4.5.2}) est vérifiée, donc $\mathcal{L}$ est ample.

Prouvons maintenant que si $q$ est de type fini et si $\mathcal{L}$ est
ample, il vérifie la condition e). En premier lieu, il résulte du critère b)
de (\hyperref[II.4.5.2-fr]{4.5.2}) et de
(\hyperref[II.4.4.1-fr]{4.4.1, (i)}) qu'il existe

